% Generated draft: requires mathematical and source review.
\documentclass[11pt]{article}
\usepackage[margin=25mm]{geometry}
\usepackage{fontspec}
\setmainfont{Latin Modern Roman}
\usepackage{amsmath,amssymb,mathtools,unicode-math}
\setmathfont{Latin Modern Math}
\usepackage{enumitem,needspace,xurl,hyperref}
\setlength{\parindent}{0pt}
\setlength{\parskip}{5pt}
\setlength{\emergencystretch}{3em}
\newcommand{\N}{\mathbb{N}}
\newcommand{\Z}{\mathbb{Z}}
\newcommand{\Q}{\mathbb{Q}}
\newcommand{\R}{\mathbb{R}}
\newcommand{\C}{\mathbb{C}}
\newcommand{\F}{\mathbb{F}}
\newcommand{\A}{\mathbb{A}}
\newcommand{\PP}{\mathbb{P}}
\newcommand{\E}{\mathbb{E}}
\newcommand{\eps}{\varepsilon}
\newcommand{\dd}{\,\mathrm{d}}
\begin{document}
\section*{UnsolvedMath ID: 700644}
\textbf{Source problem code:} catalogue entry
\subsection{Definitions and mathematical statement}
Let \(G\) be a finite group. An involution is an element of order exactly \(2\), and a \(2\)-subgroup is a subgroup whose order is a power of \(2\), including the trivial subgroup. For \(X\subseteq G\), write \(\langle X\rangle\) for the subgroup generated by \(X\).

The source asks: describe the finite groups \(G\) containing a set \(D\) of involutions such that
\[
\text{for every }D_0\subseteq D,\qquad
\langle D_0\rangle\text{ is a }2\text{-subgroup}
\ \Longrightarrow\ 
\langle N_D(D_0)\rangle\text{ is a }2\text{-subgroup}.
\]
Here the notation \(N_D(D_0)\) requires clarification: \(D\) is specified only as a set, and the supplied record does not define its normalizer notation. One possible interpretation is the set of elements of \(D\) normalizing \(D_0\) setwise,
\[
N_D(D_0)=\{d\in D:d^{-1}D_0d=D_0\}.
\]
Another possible interpretation is \(D\cap N_G(\langle D_0\rangle)\), where
\[
N_G(H)=\{g\in G:g^{-1}Hg=H\}
\]
for a subgroup \(H\leq G\). These interpretations are not identified with one another here.

The source does not require \(D\) to be nonempty, to generate \(G\), or to be invariant under conjugation. It also does not exclude \(D_0=\varnothing\). Under either interpretation above, \(N_D(\varnothing)=D\), so the empty-subset case itself requires \(\langle D\rangle\) to be a \(2\)-subgroup. These unresolved scope and notation issues must be checked before treating this as a definitive classification problem; no additional hypotheses are imposed in this draft.
\subsection{Short English statement}
Describe the finite groups admitting a set of involutions with the stated property: every subset generating a \(2\)-subgroup has a normalizer within that set which also generates a \(2\)-subgroup. The meaning of this normalizer and any intended restrictions on the set or its subsets remain unresolved.
\subsection{Sources}
\begin{itemize}
\item Local UnsolvedMath registry, permanent ID \(700644\), exact source code \texttt{catalogue entry}, title: Kourovka Notebook Problem 10.29. The supplied registry labels the problem \texttt{open}; this label is not an independently verified mathematical or literature assessment.
\item The supplied attribution identifies \emph{The Kourovka Notebook (Unsolved Problems in Group Theory)}, original problem 10.29, 10th issue (1986), proposed by A.\ A.\ Makhn\"ev. It also identifies the 21st issue edition (2026), arXiv:\texttt{1401.0300v48}. The collection description names E.\ I.\ Khukhro and V.\ D.\ Mazurov as editors.
\item The supplied registry record was created and updated on \texttt{2026-10-16T14:00:00Z} and records upstream revision \texttt{d371ee917ac9647fbb06cd6fc529deec6f3dda82}. These are registry provenance details, not dates of independent source verification. The cited notebook text has not been independently consulted for this draft.
\end{itemize}
\end{document}
